\section{Постановка задачи}
\label{sec:task}

Целью лабораторной работы является изучение структуры и принципов работы баз знаний интеллектуальных систем, освоение навыков формализации предметных областей и фактографических высказываний в базах знаний ostis-систем, а также приобретение умений навигации по базе знаний.

Индивидуальный вариант: \textbf{15) Музыка}.

Согласно методическому пособию требуется:
\begin{enumerate}
	\item выбрать тему предметной области (уникальную в потоке и согласованную с преподавателем);
	\item формализовать структуру предметной области (классы объектов исследования и исследуемые отношения);
	\item описать классы объектов исследования и исследуемые отношения (суммарно не менее 30 абсолютных и относительных понятий) с теоретико-множественными связями, свойствами и доменами отношений, а также с примечаниями, пояснениями и определениями;
	\item для каждого понятия описать не менее одного экземпляра и связать экземпляры относительными понятиями по открытым источникам;
	\item протестировать фрагменты базы знаний в системе NIKA;
	\item составить электронный отчёт (цель, задание, вариант, иллюстрации тестирования, выводы);
	\item уметь отвечать на контрольные вопросы и решать задачи;
	\item опубликовать изменения в форк проекта NIKA на GitHub.
\end{enumerate}
